\section{The Sigma transition law}
\label{sec:method}

\begin{figure*}[t]
\centering
\resizebox{\textwidth}{!}{%
\begin{tikzpicture}[
  font=\small, >=Latex, node distance=4mm and 6mm,
  box/.style={draw=sigmaBlue!80, rounded corners=2pt, fill=lightBlue, align=center,
              minimum height=8mm, minimum width=19mm},
  op/.style={draw=sigmaTeal!80, rounded corners=2pt, fill=lightTeal, align=center,
             minimum height=7mm, minimum width=18mm, font=\footnotesize},
  ctrl/.style={draw=sigmaOrange!90, rounded corners=2pt, fill=lightOrange, align=center,
               minimum height=8mm, minimum width=24mm},
  state/.style={draw=black!45, dashed, rounded corners=2pt, fill=white, align=center,
                minimum height=9mm, minimum width=32mm, font=\footnotesize},
  arrow/.style={->, thick, draw=black!70},
  darrow/.style={->, thick, draw=sigmaTeal!80},
  carrow/.style={->, thick, draw=sigmaOrange!90},
  statearrow/.style={<->, dashed, draw=black!55}
]
\node[box] (tok) {token ids};
\node[box, right=of tok] (embed) {embedding\\+ residual $x$};
\node[box, right=of embed, fill=lightOrange] (norm) {$u=\rms(x)$};
\node[op, above right=10mm and 12mm of norm] (conv) {ShortConv};
\node[op, right=of conv] (delta) {delta rule};
\node[op, below=3mm of delta] (slot) {slot memory};
\node[op, below=3mm of slot] (summ) {summary ring};
\node[op, below=3mm of summ] (win) {window (1/4)};
\node[ctrl, right=18mm of slot] (ctrl) {controller\\independent sigmoid gates\\$\Delta t(u)$};
\node[ctrl, right=of ctrl] (mix) {gate-normalised mixture\\$m-\mathrm{grad}\,E(u)+u$};
\node[box, right=of mix] (proj) {$v=\Proj(\cdot)$};
\node[box, right=of proj] (out) {$W_o v$};
\node[box, right=of out] (add) {$x'=x+W_o v$};
\node[state, below=13mm of summ] (state) {carried state: conv tail, delta matrices, slot prototypes and accumulators,\\
summary ring and partial chunk, window ring, previous normalised token};

\draw[arrow] (tok) -- (embed);
\draw[arrow] (embed) -- (norm);
\foreach \op in {conv,delta,slot,summ,win} {
  \draw[darrow] (norm.east) -| (\op.west);
  \draw[darrow] (\op.east) -| (mix.west);
  \draw[statearrow] (state.north) -| (\op.south);
}
\draw[carrow] (norm.east) |- (ctrl.west);
\draw[carrow] (ctrl) -- (mix);
\draw[arrow] (mix) -- (proj);
\draw[arrow] (proj) -- (out);
\draw[arrow] (out) -- (add);
\draw[arrow] (embed.east) |- (add.south);

\node[draw=sigmaOrange!60, dashed, rounded corners, fit=(conv)(delta)(slot)(summ)(win), inner sep=2.5mm,
      label={[sigmaOrange, font=\footnotesize, yshift=2mm]above left:parallel operator bank}] {};
\end{tikzpicture}}
\caption{Implemented SigmaStorm dataflow. The normalised latent $u$ is sent to the enabled
operators in parallel; each operator reads and updates only its own carried state. The controller
produces independent gates and the step size, the gate-normalised mixture is combined with the
energy gradient and latent residual, and the projected proposal is added to the residual stream.
The diagram is a dataflow contract, not a literal tensor-layout drawing.}
\label{fig:architecture}
\end{figure*}

\subsection{The block}

A sequence is processed as a stream. For each token position $t$ the model keeps a residual
stream $x_t\in\R^{d}$ and a fixed-size recurrent state $S_t$. One block applies
\Cref{eq:block} once: normalise, offer the normalised latent to the enabled operators, mix
their proposals, subtract an energy gradient, and project. Stacking $L$ blocks applies the law
$L$ times to the token-local latent, while $S_t$ carries information across tokens.

Three details are not cosmetic.
\begin{itemize}[leftmargin=1.1em]
\item \textbf{Independent sigmoid gates.} A softmax over operators forces them to compete for
      a fixed total and can suppress an operator that is useful on a particular token. Independent
      sigmoids let any subset contribute simultaneously; they are mixture weights, not a compute
      skipping mechanism.
\item \textbf{Convex normalisation.} Dividing by the gate sum keeps the proposal scale from
      growing merely because more operators are enabled. The normalisation is a stability device,
      not evidence that any particular routing distribution is optimal.
\item \textbf{Per-operator RMSNorm.} Normalising each proposal before gating puts heterogeneous
      operator outputs on a comparable scale, so the gate is not dominated by an arbitrary raw
      output magnitude.
\end{itemize}

The step size is $\Delta t(u)=\Delta t\,(0.5+\sigma(a^\top u))$ with $\Delta t=0.5$ by default.
The projection acts on the latent $v$ only, never on the residual stream; this is what keeps
token identity available to the associative operator.

\subsection{Operator family}

\Cref{tab:operators} lists the operators, their role, intended fixed-configuration cost, and
state size. Window attention is enabled on one block in four by default and can be removed
entirely; it retains a bounded ring history. The current implementation also materialises a mask
over each input chunk, so the strict bounded-window runtime expression below should not be read as
a guarantee for arbitrarily large unchunked calls.

\begin{table*}[t]
\centering
\caption{Operators in the SigmaStorm block. Costs are per sequence of length $T$; state is
per sample and independent of $T$. $k$ is the convolution kernel, $H$ the number of delta
heads, $k_{\mathrm{exp}}$ the state-expansion factor, $d_h=d/H$ the head width, $S$ the number
of summary slots, and $W$ the window length.}
\label{tab:operators}
\small
\begin{tabularx}{\textwidth}{@{}l X l l@{}}
\toprule
Operator & Role & Cost & State \\
\midrule
ShortConv & local order; the fast proposal & $O(Tkd)$ & $k-1$ tokens \\
Gated delta rule & associative recall (key$\to$value bindings) & linear in $T$ for fixed chunks and width & $H k_{\mathrm{exp}}$ matrices $d_h\times d_h$ \\
Slot memory & long-horizon prototypes with per-token gated writes & $O(T n_s d)$ & $n_s\times d$ + accumulators \\
Chunk summaries & bounded global summary path without $T^2$ & $O(TSd)$ & ring of $S$ summaries \\
Window attention & precise local copying (1 block in 4, optional) & bounded-window intent; chunk-mask caveat above & ring buffer $W\times d$ \\
Energy potential & contraction and stabiliser, $-\nabla E$ & $O(Td r)$ for rank $r$ & none \\
Bounded projection & hard norm bound $\|v\|\leq R$ & $O(Td)$ & none \\
Gated FFN / MoE & per-token compute & $O(Td\,d_{\mathrm{ff}})$ & none \\
\bottomrule
\end{tabularx}
\end{table*}

\subsection{Energy contraction}
\label{sec:method-energy}

The potential is
\begin{equation}
  \begin{aligned}
  E(u)&=\tfrac12 u^\top D u+\tfrac12\|L^\top u\|^2,\\
  -\nabla E(u)&=-(Du+LL^\top u),
  \end{aligned}
\end{equation}
with $D$ diagonal positive and $L\in\R^{d\times r}$ low-rank. $E$ is strongly convex, so
$-\nabla E$ is an anisotropic contraction towards the origin. It is a stabiliser, not a
multi-basin memory: strong convexity forbids many equilibria by construction. The slot memory
and the delta rule are what hold information; the energy term keeps the latent from drifting
and makes the projected dynamics well behaved. An earlier design that treated the latent as a
multi-well attractor produced no usable memory and was discarded.

\subsection{Bounded projection}

The projection is $\Proj(z)=z\cdot\min(1,R/\|z\|)$ with $R=16$. It is applied after the
transition and before the output projection. Because the residual stream is additive
($x'=x+W_ov$), the bound constrains the \emph{proposal}, not the accumulated representation;
this distinction was one of the three fixes that took MQAR from chance to $0.97$.

\subsection{State contract and chunk invariance}

\begin{observation}[Carried-state chunking contract]
\label{obs:chunk}
For a fixed configuration, carrying the state $S_t$ of \Cref{tab:state} between contiguous
chunks is intended to make the hidden representation independent of the caller's chunk
boundaries. The current gate tests one split at position 50, inside both a 64-token slot block
and a 32-token summary block; it reports a maximum hidden-state difference of
$1.6\times10^{-6}$. This is an implementation check for that configuration, not an empirical
claim over every possible boundary, length, or hardware backend.
\end{observation}

\paragraph{Why the contract is plausible.}
Every operator in \Cref{tab:operators} is implemented as a causal scan whose value at position
$t$ depends on its carried state and tokens up to $t$. The short convolution carries its $k-1$
token tail; the delta rule carries its matrices and end-of-chunk state; the slot memory carries
the ring of prototypes \emph{and} unconsolidated write accumulators; the summary ring carries
completed summaries and its partial-chunk accumulator; and the window carries its ring buffer
and length. The controller and energy are token-local. These are the design conditions needed
for boundary independence. Gradient detaching changes the backward graph, not the forward
state update.

\begin{table}[t]
\centering
\caption{Carried state per block. No field has a dimension that grows with the processed length
$T$ for a fixed configuration. The slot and summary accumulators preserve partial work across
caller chunks; the reported equivalence result is limited to the tested configuration.}
\label{tab:state}
\small
\begin{tabular}{@{}l l@{}}
\toprule
Field & Shape / content \\
\midrule
\texttt{conv\_tail} & $[B,k-1,d]$ \\
\texttt{ssm} & $[B,Hk_{\mathrm{exp}},d_h,d_h]$ \\
\texttt{memory} & keys, values, usage, write accumulators \\
\texttt{summary} & $[B,H,S,d_h]^{\times 2}$ + usage/pos + partial chunk \\
\texttt{window} & $[B,H,W,d_h]^{\times 2}$ \\
\texttt{win\_len} & scalar \\
\texttt{prev} & $[B,d]$ (previous normalised token) \\
\bottomrule
\end{tabular}
\end{table}

\subsection{Streaming forward and update lifecycle}
\label{sec:algorithm}

\begin{figure*}[t]
\centering
\fbox{\begin{minipage}{0.94\textwidth}
\textbf{Algorithm 1: carried-state training step.}
\begin{enumerate}[leftmargin=1.6em]
\item Initialise the block states once at the beginning of a sequence. Do not reset them at a
      caller-defined micro-chunk boundary.
\item For each micro-chunk, run every block from the incoming state. Within a block, compute
      $u=\rms(x)$, evaluate the enabled operators from $u$ and their carried states, form the
      gate-normalised mixture, apply the energy step and bounded projection, then add the
      projected proposal to the residual stream. Return the updated state with the chunk output.
\item Accumulate language-model and configured multi-token-prediction losses from flat token
      slices. Repeat for \texttt{bptt\_chunks} micro-chunks while carrying the forward state.
\item Detach the carried state at the truncated-backpropagation boundary, backpropagate the
      accumulated loss, unscale and clip gradients when using mixed precision, then perform the
      optimiser update and clear the accumulation window.
\item Save checkpoints at the configured cadence, including the architecture configuration and
      optimiser state needed to resume. Evaluation and generation carry state across chunks but
      do not perform optimiser updates.
\end{enumerate}
\end{minipage}}
\caption{The implemented streaming lifecycle. State detachment limits gradient history; it does
not reset the forward memory. The algorithm describes the current training/evaluation interface,
not an instruction-tuning or preference-optimisation pipeline.}
\label{fig:stream-algorithm}
\end{figure*}

\subsection{Training interface}

Training processes \texttt{micro\_chunk} tokens per forward call (default 512 in the current
training interface) and carries the state within a sequence. The differentiable horizon is
$\texttt{bptt\_chunks}\times\texttt{micro\_chunk}$ (default $4\times512=2048$ tokens); the
state is detached at the window boundary. One reported Kaggle training profile used
256-token micro-chunks, so the paper reports the configured value rather than treating one
profile as a universal default. Forward context is not limited by this horizon, but gradients
are. Distributed execution uses DDP with \texttt{no\_sync} on all but the final backward and
an fp16/bf16 scaler protocol of unscale--clip--step--update. The schedule is warmup--stable--decay;
multi-token-prediction heads are ramped from zero; the sequence length can be ramped from short
windows to long ones. Cross-entropy is computed over flat token slices, never as a
$[B,T,\text{vocab}]$ tensor. Full hyperparameters are in \Cref{app:configs}; the defaults are
interface defaults, not a claim that every reported run used them.

\subsection{Implementation invariants}

Five implementation invariants were found while bringing the stack from chance-level to working
recall. They are documented because violating any one can silently damage the result:
\begin{enumerate}[leftmargin=1.2em]
\item \textbf{The delta-rule key is the previous token.} A binding is written when the value
      arrives at $t$; its key is the token at $t-1$. Reading query, key, and value all from
      $t$ forces the projections to disentangle a blended pair, and they do not.
\item \textbf{The delta rule reads the clean stream, and the decay bias is protected from the
      generic initialiser.} A generic ``initialise every bias to zero'' pass resets
      $\alpha=\exp(-\mathrm{softplus}(\cdot))$ to $0.5$ per token, and bindings vanish in two
      steps.
\item \textbf{The transition is an additive residual with a convex mixture.} Replacing the
      residual stream with the raw latent saturates the bound and erases token identity.
\item \textbf{Operator scale is controlled before mixing.} Per-operator RMSNorm and
      gate-sum normalisation keep the mixture from being dominated by raw output scale or by
      the number of open gates; this is a stability condition, not a measured quality gain.
\item \textbf{MQAR keys must be unique per sequence.} Sampling with replacement creates
      ambiguous queries and caps the achievable score near $0.6$.
\end{enumerate}
